\documentclass[11pt,reqno]{amsart}

% ================================================================
% Packages
% ================================================================
\usepackage{amsmath,amssymb,amsthm}
\usepackage{mathrsfs}
\usepackage[shortlabels]{enumitem}
\usepackage{booktabs}
\usepackage{array}
\usepackage{microtype}
\usepackage[dvipsnames]{xcolor}
\usepackage[colorlinks=true,
 linkcolor=blue!60!black,
 citecolor=green!40!black,
 urlcolor=blue!60!black]{hyperref}

% ================================================================
% Theorem environments
% ================================================================
\newtheorem{theorem}{Theorem}[section]
\newtheorem{proposition}[theorem]{Proposition}
\newtheorem{lemma}[theorem]{Lemma}
\newtheorem{corollary}[theorem]{Corollary}
\newtheorem{conjecture}[theorem]{Conjecture}
\theoremstyle{definition}
\newtheorem{definition}[theorem]{Definition}
\newtheorem{construction}[theorem]{Construction}
\newtheorem{example}[theorem]{Example}
\theoremstyle{remark}
\newtheorem{remark}[theorem]{Remark}

% ================================================================
% Macros
% ================================================================
\newcommand{\cA}{\mathcal{A}}
\newcommand{\cD}{\mathcal{D}}
\newcommand{\cH}{\mathcal{H}}
\newcommand{\cM}{\mathcal{M}}
\newcommand{\barB}{\overline{B}}
\newcommand{\Ran}{\mathrm{Ran}}
\newcommand{\FM}{\overline{\mathrm{FM}}}
\newcommand{\Conf}{\mathrm{Conf}}
\newcommand{\fg}{\mathfrak{g}}
\newcommand{\Sym}{\mathrm{Sym}}
\newcommand{\Res}{\operatorname{Res}}
\newcommand{\Spec}{\operatorname{Spec}}
\newcommand{\Diff}{\mathrm{Diff}}
\newcommand{\Vir}{\mathrm{Vir}}
\newcommand{\Hom}{\mathrm{Hom}}
\newcommand{\id}{\mathrm{id}}
\renewcommand{\Bbbk}{\mathbf{C}}

\DeclareMathOperator{\gr}{gr}

\numberwithin{equation}{section}

% ================================================================
\begin{document}

\title[Modular Koszul Duality: a Programme]
{Modular Koszul Duality for Chiral Algebras\\
on Algebraic Curves: a Programme}

\author{Raeez Lorgat}
\address{Perimeter Institute for Theoretical Physics,
 Waterloo, ON N2L 2Y5, Canada}
\email{rlorgat@perimeterinstitute.ca}

\date{\today}

\begin{abstract}
Modular Koszul duality for chiral algebras is $E_1$--$E_1$
operadic Koszul duality transported into the homotopical
modular chiral realm on algebraic curves.
\end{abstract}

\maketitle

% ====================================================================
\section{The bar construction on algebraic curves}
\label{sec:bar}
% ====================================================================

The bar construction of an associative algebra is classical:
given an augmented algebra~$A$ over a field, form the tensor
coalgebra $T^c(s^{-1}\bar A) = \bigoplus_{n \ge 1}
(s^{-1}\bar A)^{\otimes n}$ on desuspended generators, equip it
with the differential that encodes the multiplication, and
observe that $d^2 = 0$ is equivalent to associativity.
Koszul duality, homological algebra, and deformation theory
all descend from this single construction.

A chiral algebra on a smooth algebraic curve~$X$ encodes
operator product singularities: the product of two sections
$a(z)\,b(w)$ diverges as $z \to w$, and the divergence carries
algebraic data. The problem is to build the bar construction in
this setting. The point is replaced by a curve; the
multiplication map is replaced by a chiral bracket with poles
along the diagonal; the tensor coalgebra must account for the
geometry of colliding points on~$X$.

% ====================================================================
\subsection{The propagator}
\label{ssec:propagator}
% ====================================================================

On the configuration space $C_2(X) = X \times X \setminus \Delta$
of two distinct points on a smooth algebraic curve~$X$
over~$\Bbbk$, define
\begin{equation}\label{eq:eta}
\eta_{12}
\;=\;
d\log(z_1 - z_2)
\;=\;
\frac{dz_1 - dz_2}{z_1 - z_2}\,.
\end{equation}
This $1$-form has a simple pole along the diagonal
$\Delta \subset X \times X$ with residue~$1$. The residue is
intrinsic: it depends on no coordinate, no metric, no level.

Three properties single out~$\eta_{12}$. A double pole
$dz/(z-w)^2$ transforms under coordinate change and has no
coordinate-independent residue. A regular $1$-form has
residue zero and extracts nothing from the collision. The
logarithmic derivative is the unique $1$-form on~$C_2(X)$
with a first-order pole along~$\Delta$ and conformally
invariant residue.

On $C_3(X)$, the three pullbacks
$\eta_{12}, \eta_{23}, \eta_{31}$ satisfy
\begin{equation}\label{eq:arnold}
\eta_{12} \wedge \eta_{23}
\;+\;
\eta_{23} \wedge \eta_{31}
\;+\;
\eta_{31} \wedge \eta_{12}
\;=\; 0\,.
\end{equation}
This is the \emph{Arnold relation}. It restates the
partial-fractions identity
\[
\frac{1}{(z_1 {-} z_2)(z_2 {-} z_3)}
+
\frac{1}{(z_2 {-} z_3)(z_3 {-} z_1)}
+
\frac{1}{(z_3 {-} z_1)(z_1 {-} z_2)}
= 0\,,
\]
which holds because the left side, viewed as a rational function
of~$z_1$ with $z_2, z_3$ fixed, has simple poles at $z_2$
and~$z_3$ whose residues sum to zero.
Arnold~\cite{Arnold69} proved that~\eqref{eq:arnold}, replicated
across all triples $\{i,j,k\} \subset \{1,\ldots,n\}$, generates
the complete ideal of relations in
$H^*(\Conf_n(\Bbbk), \Bbbk)$.

Under a coordinate change $z \mapsto w(z)$, the propagator
transforms as
$\eta_{ij} \mapsto \eta_{ij}
+ d\log\frac{w(z_i) - w(z_j)}{z_i - z_j}$.
Near the diagonal the correction becomes $d\log w'(z_i)$: the
standard cocycle for~$\Diff(X)$. The propagator is a BRST ghost
for diffeomorphisms; its failure to be coordinate-invariant is
the failure that the Virasoro algebra measures.


% ====================================================================
\subsection{Chiral algebras}
\label{ssec:chiral}
% ====================================================================

A \emph{chiral algebra}~$\cA$ on~$X$ is a
$\cD_X$-module equipped with a chiral bracket
\[
\mu^{\mathrm{ch}}
\;\colon\;
j_* j^*(\cA \boxtimes \cA)
\;\longrightarrow\;
\Delta_!(\cA)\,,
\]
where $j \colon X \times X \setminus \Delta
\hookrightarrow X \times X$ is the open complement and
$\Delta_! \cA$ is the $!$-pushforward along the diagonal
embedding~\cite{BD04}. The source consists of sections of
$\cA \boxtimes \cA$ with poles of arbitrary order
along~$\Delta$; the chiral bracket extracts the singular
part and projects it onto the diagonal. In local coordinates,
the chiral bracket encodes the operator product expansion:
\[
a(z)\,b(w)
\;\sim\;
\sum_{n \ge 0}
\frac{a_{(n)}b(w)}{(z - w)^{n+1}}\,.
\]
The $n$-th products $a_{(n)}b \in \cA$ are the Fourier
modes, extracted by
$a_{(n)}b
= \Res_{z=w}[a(z)\,b(w)\,(z{-}w)^n]$.
The zeroth product $a_{(0)}b$ is the Lie bracket
(first-order pole); $a_{(1)}b$ carries the
symmetric/metric data (second-order pole). Higher products
$a_{(n)}b$ for $n \ge 2$ are determined by the
$\cD_X$-module structure.

The Borcherds identity
\begin{equation}\label{eq:borcherds}
\sum_{j \ge 0} \binom{m}{j}
(a_{(n+j)}b)_{(m+k-j)}c
\;=\;
\sum_{j \ge 0} (-1)^j \binom{n}{j}
\Bigl(
a_{(m+n-j)}(b_{(k+j)}c)
- (-1)^n b_{(n+k-j)}(a_{(m+j)}c)
\Bigr)
\end{equation}
is the complete axiom system. At $n = 0$ it reduces to the
Jacobi identity for the Lie bracket~$a_{(0)}$; at $n = 1$,
$m = k = 0$, to the derivation property of~$a_{(0)}$ with
respect to~$a_{(1)}b$.


% ====================================================================
\subsection{The bar complex}
\label{ssec:bar-complex}
% ====================================================================

Place sections $a_1, \ldots, a_n \in \cA$ at distinct points
$z_1, \ldots, z_n \in X$. The Fulton--MacPherson
compactification $\FM_n(X)$ resolves $\Conf_n(X)$ by
iterated real oriented blowup along all diagonal
strata~\cite{FM94}. Points of
$\FM_n(X) \setminus \Conf_n(X)$ record the limiting direction
and relative scale of approach as points collide. The boundary
is a manifold with corners; codimension-one faces $D_S$
are indexed by subsets $S \subseteq \{1,\ldots,n\}$ with
$|S| \ge 2$, recording which points have collided. Each face
factors:
\begin{equation}\label{eq:fm-boundary}
D_S
\;\simeq\;
\FM_{|S|} \times \FM_{n - |S| + 1}\,,
\end{equation}
cluster times residual.

The \emph{ordered bar complex} is
\begin{equation}\label{eq:bar-ord}
\barB^{\mathrm{ord}}_X(\cA)
\;=\;
\bigoplus_{n \ge 1}
\bigl(s^{-1}\bar\cA\bigr)^{\otimes n}
\;\otimes\;
\Gamma\bigl(
\FM_n(X),\;
\Omega^{n-1}_{\log}
\bigr),
\end{equation}
with no $\Sigma_n$-quotient. Here
$\bar\cA = \ker(\varepsilon \colon \cA \to \omega_X)$
is the augmentation ideal, $s^{-1}$ denotes desuspension
($|s^{-1}a| = |a| - 1$), and the cohomological convention
$|d| = +1$ is used throughout.

The bar differential decomposes:
\begin{equation}\label{eq:bar-diff}
d_{\barB}
\;=\;
d_{\mathrm{int}} + d_{\mathrm{res}} + d_{\mathrm{dR}}\,.
\end{equation}
The internal differential $d_{\mathrm{int}}$ acts on each
tensor factor; the de~Rham differential $d_{\mathrm{dR}}$ acts
on logarithmic forms on~$\FM_n(X)$; and the residue differential
$d_{\mathrm{res}}$ extracts OPE data at collision divisors.
For a codimension-one face $D_{\{i,j\}}$ corresponding to the
collision $z_i \to z_j$, the residue differential extracts
the $p$-th Fourier mode $a_{(p)}b$ of the OPE, where $p$ is
the pole order read from the form-theoretic coefficient of
$\eta_{ij}$. The bar differential does not merely detect
singularities; it extracts them, mode by mode, through the
logarithmic kernel~\eqref{eq:eta}.


% ====================================================================
\subsection{Why $d^2 = 0$}
\label{ssec:d-squared}
% ====================================================================

Expand $(d_{\mathrm{int}} + d_{\mathrm{res}}
+ d_{\mathrm{dR}})^2$ into nine terms. Three vanish
individually: $d_{\mathrm{int}}^2 = 0$ (the underlying cochain
complex), $d_{\mathrm{dR}}^2 = 0$ (de~Rham), and
$d_{\mathrm{int}}\,d_{\mathrm{dR}}
+ d_{\mathrm{dR}}\,d_{\mathrm{int}} = 0$ (Leibniz). The
remaining terms cancel pairwise, governed by index overlap:

\begin{enumerate}[(a)]
\item \emph{Disjoint supports.}
When $\{i,j\} \cap \{k,l\} = \emptyset$, the faces
$D_{\{i,j\}}$ and $D_{\{k,l\}}$ are transverse
in~$\FM_n(X)$. The residue operators commute:
$d_{\mathrm{res}}^{(ij)}\,d_{\mathrm{res}}^{(kl)}
+ d_{\mathrm{res}}^{(kl)}\,d_{\mathrm{res}}^{(ij)} = 0$.

\item \emph{Shared index.}
When $\{i,j\} \cap \{k,l\} = \{j = k\}$, the triple
collision $z_i = z_j = z_l$ sits in a codimension-two
stratum. The form-theoretic contribution involves
$\eta_{ij} \wedge \eta_{jl}$. The Arnold
relation~\eqref{eq:arnold} identifies this with a cyclic
sum over the three ordered approaches to the triple-collision
point. On the algebraic side, the three terms produce
$(a_{i(m)}a_j)_{(n)}a_l$,
$(a_{j(m)}a_l)_{(n)}a_i$,
$(a_{l(m)}a_i)_{(n)}a_j$
with appropriate pole orders, whose sum vanishes by the
Borcherds identity~\eqref{eq:borcherds}.

The Arnold relation kills the form; the Borcherds identity
kills the algebra. Both are needed, and both are consumed.

\item \emph{Same pair.}
$d_{\mathrm{res}}^{(ij)}\,d_{\mathrm{res}}^{(ij)} = 0$
by the sign rule on desuspensions: extracting a residue
twice at the same diagonal vanishes for degree reasons.
\end{enumerate}

The result: $d_{\barB}^2 = 0$ if and only if the
$n$-th products satisfy~\eqref{eq:borcherds}. The bar
construction accepts exactly chiral algebras.


% ====================================================================
\subsection{Two coproducts}
\label{ssec:coproducts}
% ====================================================================

The ordered bar $\barB^{\mathrm{ord}}(\cA)$ carries the
\emph{deconcatenation coproduct}:
\begin{equation}\label{eq:deconc}
\Delta[s^{-1}a_1 | \cdots | s^{-1}a_n]
\;=\;
\sum_{i=0}^{n}
[s^{-1}a_1 | \cdots | s^{-1}a_i]
\;\otimes\;
[s^{-1}a_{i+1} | \cdots | s^{-1}a_n]\,,
\end{equation}
splitting an ordered sequence at every consecutive position:
$n + 1$ terms. This is coassociative but not cocommutative.
It makes $\barB^{\mathrm{ord}}(\cA)$ a cofree tensor coalgebra
$T^c(s^{-1}\bar\cA)$: the $E_1$-coalgebra reflecting the
linear ordering of points.

The $\Sigma_n$-coinvariant quotient
\begin{equation}\label{eq:bar-sym}
\barB^{\Sigma}_X(\cA)
\;=\;
\bigoplus_{n \ge 1}
\bigl(s^{-1}\bar\cA\bigr)^{\otimes n}_{\Sigma_n}
\;\otimes\;
\Gamma\bigl(
\FM_n(X),\;
\Omega^{n-1}_{\log}
\bigr)
\end{equation}
carries the \emph{coshuffle coproduct}, summing over all $2^n$
bipartitions of an unordered set: coassociative and
cocommutative. This is the factorization coalgebra
of Beilinson--Drinfeld on~$\Ran(X)$~\cite{BD04}, the
$E_\infty$-coalgebra.

The ordered bar is the natural primitive. The symmetric bar
is a quotient that loses data: the non-cocommutative structure
of the ordered bar carries the spectral $R$-matrix, the
Knizhnik--Zamolodchikov associator, and the full quantum group
deformation, none of which survive the passage to
$\Sigma_n$-coinvariants.

\begin{remark}[Six bar complexes]\label{rem:six-bars}
Six bar-type objects coexist. The Francis--Gaitsgory
bar~\cite{FG12} $B^{\mathrm{Lie}}(\cA)$ is the cofree
coLie coalgebra $\mathrm{Lie}^c(s^{-1}\bar\cA)$, capturing
zeroth-pole data; it embeds into~$\barB^\Sigma$ via the
inclusion $\mathrm{Lie}^c \hookrightarrow \Sym^c$
(a quasi-isomorphism in characteristic zero, but a categorical
distinction). The genus-graded bar $B^{(g)}(\cA)$ extends
$\barB^\Sigma$ to families of curves over
$\overline{\cM}_g$, where the differential acquires curvature
$d_{\mathrm{fib}}^2 = \kappa(\cA) \cdot \omega_g$
(see~\S\ref{ssec:genus}).
The modular bar $B^{\mathrm{mod}}(\cA)$ restores $D^2 = 0$
by completing over the Feynman transform of the commutative
modular operad. The $E_1$-modular bar
$B^{E_1\text{-}\mathrm{mod}}(\cA)$, an algebra over the
Feynman transform of the associative modular operad, is the
primitive among these; the modular bar is its
$\Sigma_n$-coinvariant image. No two of these six objects
are isomorphic in general.
\end{remark}


% ====================================================================
\subsection{Three operations on the bar}
\label{ssec:three-operations}
% ====================================================================

The bar construction is a categorical logarithm:
$B(\cA \otimes \cA') \simeq B(\cA) \oplus B(\cA')$.
The logarithmic form
$\eta_{ij} = d\log(z_i - z_j)$ is the integral kernel;
the Arnold relation is the cocycle condition for~$\log$.

Three distinct functors act on~$B(\cA)$ and produce three
distinct outputs:
\begin{equation}\label{eq:three-functors}
\begin{array}{lcll}
\Omega(B(\cA)) &\simeq& \cA
 & \text{(cobar inversion: recovers the original)},\\[3pt]
\mathbb{D}_{\Ran}(B(\cA)) &\simeq& \cA^!_\infty
 & \text{(Verdier duality: the homotopy Koszul dual)},
 \\[3pt]
(H^*(B(\cA)))^\vee &=& \cA^!
 & \text{(linear duality of bar cohomology: the strict dual)}.
\end{array}
\end{equation}
Cobar inversion recovers the original algebra~$\cA$ itself;
it is a round-trip, not a duality. Verdier duality on
$\Ran(X)$ converts the bar coalgebra into a factorization
algebra~$\cA^!_\infty$, the homotopy Koszul dual retaining
full chain-level data. That its cohomology recovers the
strict Koszul dual~$\cA^!$ on the Koszul locus is the content
of the Verdier intertwining theorem, not a tautology.
Confusing any two of~\eqref{eq:three-functors} is a
category error.


% ====================================================================
\subsection{The Heisenberg algebra at bar degree two}
\label{ssec:heisenberg}
% ====================================================================

The Heisenberg chiral algebra $\cH_k$ is generated by a
single field~$\alpha$ of conformal weight~$1$ with OPE
\begin{equation}\label{eq:heis-ope}
\alpha(z)\,\alpha(w)
\;\sim\;
\frac{k}{(z-w)^2}\,.
\end{equation}
Only the second-order pole is present:
$\alpha_{(1)}\alpha = k \cdot \mathbf{1}$ and
$\alpha_{(n)}\alpha = 0$ for $n \neq 1$. No first-order
pole: $\alpha_{(0)}\alpha = 0$, no Lie bracket.
All triple products vanish since
$\mathbf{1}_{(n)} = 0$ for $n \ge 0$.

The bar complex at tensor degree~$2$ gives
\begin{equation}\label{eq:heis-bar2}
d_{\mathrm{res}}
(s^{-1}\alpha \otimes s^{-1}\alpha \otimes \eta_{12})
\;=\;
\alpha_{(1)}\alpha
\;=\;
k \cdot \mathbf{1}\,.
\end{equation}
The level~$k$ is visible already at genus~$0$. At tensor
degree $n \ge 3$: every residue at a collision stratum
$D_{\{i,j\}}$ produces $\alpha_{(1)}\alpha = k \cdot \mathbf{1}$
on the pair~$(i,j)$; since $\mathbf{1}_{(p)}\alpha = 0$ for
$p \ge 0$, no further contractions occur. All structure is
captured at tensor degree~$2$: the bar complex is quadratic.
The Heisenberg algebra is chirally Koszul.

The ordered bar carries the collision residue
\begin{equation}\label{eq:heis-r}
r(z) \;=\; \frac{k}{z}\,.
\end{equation}
The second-order OPE pole $k/(z-w)^2$ drops by one order under
the logarithmic kernel $d\log(z-w)$: the bar differential
absorbs one power of $(z-w)$, producing the first-order
$r$-matrix pole. Applying the averaging map:
\begin{equation}\label{eq:heis-kappa}
\mathrm{av}\bigl(r(z)\bigr)
\;=\;
k
\;=\;
\kappa(\cH_k)\,.
\end{equation}
For the Heisenberg algebra, the $R$-matrix \emph{is}
the modular characteristic: $r(z) = k/z$ is scalar, so
averaging loses nothing.

The deficiency: no first-order pole, no Lie bracket, no
matrix structure in the collision residue, no quantum group.
Everything that makes representation theory nontrivial is
absent. The simplest algebra where the $R$-matrix is
matrix-valued is the first algebra with a first-order pole:
the affine Kac--Moody algebra.

\begin{example}[Kac--Moody collision residue]
\label{ex:km}
The affine Kac--Moody algebra $\widehat{\fg}_k$ at
level~$k$ is generated by currents $J^a$,
$a = 1,\ldots,\dim\fg$, with OPE
\[
J^a(z)\,J^b(w)
\;\sim\;
\frac{f^{ab}{}_{c}\,J^c(w)}{z - w}
\;+\;
\frac{k\,\kappa^{ab}}{(z - w)^2}\,.
\]
Both poles are present. The ordered bar carries the collision
residue
\begin{equation}\label{eq:km-r}
r(z)
\;=\;
\frac{k\,\Omega}{z}\,,
\qquad
\Omega = \sum_a J^a \otimes J_a
\;\in\;
\fg \otimes \fg\,,
\end{equation}
where $\Omega$ is the Casimir tensor and $k$ is the level.
At $k = 0$ the $r$-matrix vanishes: no level, no collision
data. The averaging map collapses $k\,\Omega$ to a scalar:
\begin{equation}\label{eq:km-kappa}
\mathrm{av}(k\,\Omega/z)
\;=\;
\frac{k\,\dim(\fg)}{2h^\vee}
\;=\;
\kappa_{\mathrm{dp}}(\widehat{\fg}_k)\,,
\end{equation}
where $h^\vee$ is the dual Coxeter number. The full modular
characteristic adds the Sugawara shift:
\[
\kappa(\widehat{\fg}_k)
\;=\;
\kappa_{\mathrm{dp}}(\widehat{\fg}_k) + \frac{\dim(\fg)}{2}
\;=\;
\frac{\dim(\fg)\,(k+h^\vee)}{2h^\vee}.
\]
The matrix
structure of~$\Omega$, erased by $\Sigma_2$-coinvariance,
encodes the representation theory of~$\fg$: it is the datum
that builds the Yangian~$Y(\fg)$.
\end{example}


% ====================================================================
\subsection{Genus~$0$ and genus $\ge 1$}
\label{ssec:genus}
% ====================================================================

Everything above takes place at genus~$0$: the propagator
$\eta_{12} = d\log(z_1 - z_2)$ is globally defined
on~$\mathbb{P}^1$, and the Arnold relation holds without
correction. On a curve of genus $g \ge 1$, the function
$z_1 - z_2$ has no global meaning: the curve is compact,
and any meromorphic function with a single simple zero must
also have a simple pole.

The replacement is the \emph{prime form} $E(z_1, z_2)$, a
section of $K^{-1/2} \boxtimes K^{-1/2}$ vanishing to first
order along the diagonal with no other zeros. The genus-$g$
propagator
\begin{equation}\label{eq:genus-g-prop}
\eta^{(g)}_{12}
\;=\;
d\log E(z_1, z_2)
+ \text{(period correction)}
\end{equation}
is single-valued but no longer holomorphic; the period
correction couples to the Hodge bundle
$\mathbb{E} \to \overline{\cM}_g$ through the period
matrix. The propagator $d\log E(z_1,z_2)$ has conformal
weight~$1$ in both variables, regardless of the conformal
weight of the field being sewed: all edges of the genus-$g$
graph sum use the standard Hodge bundle
$\mathbb{E}_1 = R^0\pi_*\omega_\pi$.

The Arnold relation acquires a defect, and the fibrewise bar
differential satisfies
\begin{equation}\label{eq:curvature}
d_{\mathrm{fib}}^{\,2}
\;=\;
\kappa(\cA) \cdot \omega_g\,,
\end{equation}
where $\omega_g$ is the Arakelov form on the fibre of
$\pi \colon \mathcal{C}_g \to \overline{\cM}_g$ and
$\kappa(\cA) \in \Bbbk$ is the \emph{modular
characteristic}, determined entirely by the leading OPE
singularity. Explicit values:
\begin{equation}\label{eq:kappa-values}
\kappa(\cH_k) = k\,,
\qquad
\kappa(\widehat{\fg}_k)
= \frac{\dim(\fg) \cdot (k + h^\vee)}{2h^\vee}\,,
\qquad
\kappa(\Vir_c)
= \frac{c}{2}\,.
\end{equation}
Period integrals restore nilpotence: the Gauss--Manin
connection absorbs the fibrewise curvature, and the total
differential satisfies $D_g^2 = 0$.

Genus~$0$ is uncurved algebra. Genus $\ge 1$ is curved
geometry. The entire genus tower is controlled by a single
scalar: the curvature of the bar complex on a family of
curves.

\begin{remark}[Recovering the classical bar]
\label{rem:classical}
When $X = \Spec\Bbbk$ is a point, configuration spaces
collapse, logarithmic forms vanish, and the chiral bar complex
recovers the classical bar construction of an associative
algebra~$A$. This recovery requires specifying the retraction
data: vertex algebras live on the formal disk~$D$, and the
passage from~$D$ to a bare point discards the thickening
(completion, growth conditions) on which the
$\mathcal{D}$-module structure depends.
A quadratic algebra $A = T(V)/(R)$ has quadratic
dual $A^! = T(V^*)/(R^\perp)$; when~$A$ is Koszul, the bar
cohomology concentrates in diagonal degrees and
$H^*(B(A)) \cong (A^!)^\vee$. The operadic instances
$\mathrm{Com}^! = \mathrm{Lie}$ and
$\Sym^! = \Lambda$ recover the Chevalley--Eilenberg complex
as a restricted bar construction.
\end{remark}


% ====================================================================
% References
% ====================================================================

\begin{thebibliography}{99}

\bibitem{Arnold69}
V.\,I.~Arnold,
The cohomology ring of the group of dyed braids,
\textit{Mat.\ Zametki}~\textbf{5} (1969), 227--231.

\bibitem{BD04}
A.~Beilinson and V.~Drinfeld,
\textit{Chiral Algebras},
AMS Colloquium Publications, vol.~51,
American Mathematical Society, Providence, RI, 2004.

\bibitem{BGS96}
A.~Beilinson, V.~Ginzburg, and W.~Soergel,
Koszul duality patterns in representation theory,
\textit{J.~Amer.\ Math.\ Soc.}~\textbf{9} (1996), 473--527.

\bibitem{CG17}
K.~Costello and O.~Gwilliam,
\textit{Factorization Algebras in Quantum Field Theory},
vols.~1--2, Cambridge University Press, 2017/2021.

\bibitem{FG12}
J.~Francis and D.~Gaitsgory,
Chiral Koszul duality,
\textit{Selecta Math.\ (N.S.)}~\textbf{18} (2012), 27--87.

\bibitem{FM94}
W.~Fulton and R.~MacPherson,
A compactification of configuration spaces,
\textit{Ann.\ of Math.\ (2)}~\textbf{139} (1994), 183--225.

\bibitem{LV12}
J.-L.~Loday and B.~Vallette,
\textit{Algebraic Operads},
Grundlehren der Mathematischen Wissenschaften, vol.~346,
Springer, Berlin, 2012.

\bibitem{Mok25}
C.-P.~Mok,
\emph{Logarithmic Fulton--MacPherson configuration spaces},
arXiv:2503.17563, 2025.

\bibitem{MS24}
N.~Markarian and H.~L.~Tanaka,
Homotopy chiral algebras,
preprint, 2024.

\bibitem{Priddy70}
S.~Priddy,
Koszul resolutions,
\textit{Trans.\ Amer.\ Math.\ Soc.}~\textbf{152}
(1970), 39--60.

\end{thebibliography}

\end{document}
